\begin{textsample}{POS Dim 3 – rn – Score 14.00 – rn/rn\_inf\_20.txt}  \label{ex:f3_pos_060}
As experiências educativas aparecem e são propostas, então, em oposição ao currículo organizado por eixos/blocos/áreas de conhecimentos e/ou disciplinas curriculares.
Oliveira ( 2015 ) ressalta que as experiências foram pensadas por possibilitarem o reconhecimento do protagonismo da \textbf{criança}, a valorização do sentido pessoal que cada \textbf{criança} empresta às vivências propostas e aos conhecimentos nelas construídos e o caráter prático-reflexivo que devem assumir as práticas pedagógicas propostas às \textbf{crianças}.
O desafio que se coloca para a elaboração curricular, segundo Oliveira ( 2013 ), consiste em transcender a prática pedagógica centrada no(a) professor(a) e garantir experiências que possibilitem “ o encontro de explicações pela \textbf{criança} sobre o que ocorre à sua \textbf{volta} e consigo mesma, enquanto desenvolvem formas de sentir, pensar e solucionar problemas ” ( Ibid. ; p. 06 ).
Deste modo, não se trata de pura transmissão de uma cultura considerada pronta, mas, de promover interações nas quais as \textbf{crianças} busquem compreender o mundo e a si mesmas, através de processos de significação e construção de cultura, de um modo muito próprio e marcado pelo momento histórico.
Na proposta do currículo organizado por experiências, Augusto ( 2013 ) esclarece que os saberes das \textbf{crianças} devem ser validados pela escola e considerados desde o planejamento do professor, visando à sua articulação aos novos conhecimentos.
Desse modo, espera -se que a \textbf{criança} possa envolver -se em processos de significação tomando os novos conhecimentos e diferentes modos de aprender como parte de sua própria experiência.
Define -se assim, que sejam garantidas experiências educativas que envolvem objetivos que visam à aprendizagem dos conhecimentos construídos socialmente em diferentes âmbitos e das diferentes linguagens em suas variadas formas de expressão e interação, no contexto de currículos e práticas cotidianas da Educação \textbf{Infantil}.
Tais experiências e sua materialização nas práticas curriculares têm como eixos norteadores as interações e a \textbf{brincadeira}, que, por sua vez, constituem os processos de aprendizagem e desenvolvimento da \textbf{criança}, conforme amplamente discutido no capítulo três deste documento.
DIALOGANDO COM AS DCNEI Art. 9º As práticas pedagógicas que compõem a proposta curricular da Educação \textbf{Infantil} devem ter como eixos norteadores as interações e a \textbf{brincadeira}, garantindo experiências que : I - promovam o conhecimento de si e do mundo por meio da ampliação de experiências \textbf{sensoriais}, expressivas, corporais que possibilitem \textbf{movimentação} ampla, expressão da individualidade e respeito pelos ritmos e \textbf{desejos} da \textbf{criança} ; II - favoreçam a imersão das \textbf{crianças} nas diferentes linguagens e o progressivo domínio por elas de vários gêneros e formas de expressão : gestual, verbal, \textbf{plástica}, \textbf{dramática} e musical ; III - possibilitem às \textbf{crianças} experiências de narrativas, de apreciação e interação com a linguagem oral e escrita, e convívio com diferentes suportes e gêneros textuais orais e escritos ; IV - recriem, em contextos significativos para as \textbf{crianças}, relações quantitativas, medidas, formas e orientações espaço temporais ; V - ampliem a confiança e a participação das \textbf{crianças} nas atividades individuais e coletivas ; VI - possibilitem situações de aprendizagem mediadas para a elaboração da autonomia das \textbf{crianças} nas ações de \textbf{cuidado} pessoal, auto-organização, saúde e bem-estar ; VII - possibilitem vivências éticas e estéticas com outras \textbf{crianças} e grupos culturais, que alarguem seus padrões de referência e de identidades no diálogo e reconhecimento da diversidade ; VIII - incentivem a \textbf{curiosidade}, a exploração, o encantamento, o questionamento, a indagação e o conhecimento das \textbf{crianças} em relação ao mundo físico e social, ao tempo e à natureza ; IX - promovam o relacionamento e a interação das \textbf{crianças} com diversificadas manifestações de música, artes \textbf{plásticas} e gráficas, cinema, fotografia, dança, teatro, poesia e literatura ; X - promovam a interação, o \textbf{cuidado}, a preservação e o conhecimento da biodiversidade e da sustentabilidade da vida na Terra, assim como o não \textbf{desperdício} dos recursos naturais ; XI - propiciem a interação e o conhecimento pelas \textbf{crianças} das manifestações e tradições culturais brasileiras ; XII - possibilitem a utilização de gravadores, projetores, computadores, máquinas \textbf{fotográficas}, e outros recursos tecnológicos e midiáticos ( BRASIL, 2009a, p. 4 ).
A Base Nacional Comum Curricular ( BRASIL, 2017 ), ancorada nas Diretrizes Curriculares Nacionais para Educação \textbf{Infantil} ( BRASIL, 2009a ), propõe o currículo da Educação \textbf{Infantil} organizado por Campos de experiências a partir de Direitos de aprendizagem e desenvolvimento da \textbf{criança}.
No capítulo que segue são apresentados os Direitos e Objetivos de Aprendizagem e Desenvolvimento a serem garantidos no desenvolvimento dos cinco Campos de Experiências definidos na BNCC.

% matched lemmas: brincadeira, criança, cuidado, curiosidade, desejo, desperdício, dramático, fotográfico, infantil, movimentação, plástico, sensorial, volta
\end{textsample}
